\documentclass[12pt,a4paper]{article}
\usepackage[utf8]{inputenc}
\usepackage[T1]{fontenc}
\usepackage{amsmath}
\usepackage{amsfonts}
\usepackage{amssymb}
\usepackage{graphicx}
\usepackage{hyperref}
\usepackage{geometry}
\usepackage{booktabs}
\usepackage{listings}
\usepackage{xcolor}
\usepackage{float}

\geometry{margin=1in}

\title{\textbf{Intelligent Solar Energy Generation Forecasting and Agentic Grid Optimization System}}
\author{Team 13: Dhiraj (TL), Shorya, Ved, Himanshu}
\date{April 2026}

\definecolor{codegreen}{rgb}{0,0.6,0}
\definecolor{codegray}{rgb}{0.5,0.5,0.5}
\definecolor{codepurple}{rgb}{0.58,0,0.82}
\definecolor{backcolour}{rgb}{0.95,0.95,0.92}

\lstset{
    backgroundcolor=\color{backcolour},   
    commentstyle=\color{codegreen},
    keywordstyle=\color{magenta},
    numberstyle=\tiny\color{codegray},
    stringstyle=\color{codepurple},
    basicstyle=\ttfamily\footnotesize,
    breakatwhitespace=false,         
    breaklines=true,                 
    captionpos=b,                    
    keepspaces=true,                 
    numbers=left,                    
    numbersep=5pt,                  
    showspaces=false,                
    showstringspaces=false,
    showtabs=false,                  
    tabsize=2
}

\begin{document}

\maketitle

\section{Introduction}
As global energy consumption shifts towards renewable sources, solar power generation has become a critical component of the sustainable energy mix. However, the inherent variability of solar energy due to weather conditions poses significant challenges for grid stability and energy management. This project develops an end-to-end platform that combines high-accuracy Machine Learning forecasting with an Agentic AI assistant to provide actionable grid optimization recommendations.

\section{System Objectives}
The primary goals of the system are:
\begin{itemize}
    \item To provide precise hour-ahead and day-ahead solar power generation forecasts based on meteorological data.
    \item To implement an intelligent AI agent that reasons about these forecasts to identify risks and suggest optimization strategies.
    \item To deliver a professional, user-centric dashboard for real-world energy management.
    \item To provide structured, automated reports for grid balancing and battery storage coordination.
\end{itemize}

\section{Data and Preprocessing}
The system utilizes a dataset of over 4,200 weather observation records. Key features include:
\begin{itemize}
    \item \textbf{Meteorological}: Temperature, Relative Humidity, Pressure, Wind Speed/Direction, and Cloud Cover.
    \item \textbf{Solar}: Shortwave Radiation and Sky Geometry (Zenith, Azimuth, Angle of Incidence).
\end{itemize}
Data preprocessing includes outlier removal, timestamp synchronization, and feature scaling through the \texttt{src/data\_preprocessing.py} pipeline.

\section{Machine Learning Engine}
The foresting engine employs a \textbf{Random Forest Regressor} ensemble model. The model was optimized through hyperparameter tuning (n\_estimators=200, max\_depth=15) and achieves the following performance:
\begin{itemize}
    \item \textbf{Coefficient of Determination ($R^2$)}: 0.9312
    \item \textbf{Mean Absolute Error (MAE)}: 187.5 kW
\end{itemize}
The high $R^2$ value indicates that 93\% of the variance in power generation is captured by the meteorological features.

\section{Agentic AI Optimization Assistant}
The core innovation of Milestone 2 is the "Ved AI" Agentic Assistant. Unlike traditional LLM chatbots, this system uses a structured reasoning graph:
\begin{enumerate}
    \item \textbf{Forecast Analyst Node}: Interprets raw prediction values and metadata.
    \item \textbf{Knowledge Retrieval (RAG)}: Uses FAISS to search a custom database of grid management guidelines and battery safety protocols.
    \item \textbf{Optimization Synthesis}: Combines data and knowledge into a structured JSON report using Google Gemini 1.5 Flash.
\end{enumerate}

\section{User Interface and Visualization}
The dashboard, built with Streamlit and Plotly, provides three core modules:
\begin{itemize}
    \item \textbf{Dashboard Analytics}: High-level metrics showing estimated savings and average performance.
    \item \textbf{Interactive Forecast}: Real-time weather input forms with 24-hour power curve simulations.
    \item \textbf{Optimization Center}: Interface for the Agentic Assistant and PDF report generation.
\end{itemize}

\section{Grid Management Recommendations}
Based on the system's reasoning, final users receive structured advice covering:
\begin{itemize}
    \item \textbf{Battery Coordination}: Charge-discharge cycles based on predicted peak radiation hours.
    \item \textbf{Risk Mitigation}: Warning of low generation periods due to cloud cover.
    \item \textbf{Load Balancing}: Suggestions for industrial consuming offsets during high generation bursts.
\end{itemize}

\section{Conclusion}
This system demonstrates the power of combining traditional ML with modern Agentic AI. By transforming sparse meteorological data into human-understandable optimization strategies, we bridge the gap between "predictive analysis" and "intelligent action," satisfying all requirements for the solar grid optimization system.

\end{document}
